\documentclass[10pt]{letter}
\usepackage[margin=0.8in]{geometry}
\usepackage{microtype}
\signature{Mat\'eo H. Petel, M.S.}
\address{Center for Strategic Research in Critical Technologies; SentinelOps\\Building 208, Rosse Lane, Stanford, CA 94305, USA\\+1 650 223 4809; mpetel@stanford.edu}
\begin{document}
\begin{letter}{Editorial Office\\Digital Discovery\\Royal Society of Chemistry}
\opening{Dear Editors,}

Please consider ``Operational Bottlenecks in the Digital Laboratory: A Reproducible Mixed-Methods Study of Public Practitioner Evidence'' as a Full Paper in Digital Discovery.

The manuscript examines an underrepresented layer of digital-laboratory science: the operational evidence required to connect device interfaces, deployable methods, physical-resource models, event records, recovery, scheduling, and stage-bounded validation. We analyze 55 purposively selected public laboratory-automation discussions and 45 episodes from a deliberately difficult subset. The study contributes a ten-construct model, bounded metrics, and robustness analyses while preserving explicit limits on prevalence, causal, and vendor-comparison inference.

The submission is also a reproducible digital research resource. The versioned release contains derived data, field-level metric inputs, source and claim audit records, JSON Schemas, executable analysis code, deterministic figures and tables, regression tests, and governed example records. Referees can reproduce every reported numerical result from committed records without accessing the forum. The Data Availability Statement identifies the public release; a persistent archival DOI will be inserted before final submission.

The paper fits Digital Discovery through its treatment of digital laboratory infrastructure, interoperable automation, research-data stewardship, AI-method evaluation, and machine-readable operational evidence. Its strongest mechanism links command--state--material observability to recovery after irreversible physical execution. The resulting schemas and prospective outcomes provide a path from the current pilot to denominated operational studies.

The manuscript is original, is not under consideration elsewhere, and has not been published as an archival article. Thank you for considering it.

\closing{Sincerely,}
\end{letter}
\end{document}
